% EdgeSched: phase-aware placement and compute sharing -- core model
% This standalone note mirrors the current system design. It deliberately
% avoids routing by model name or by a single CPU/GPU speed ratio.

\documentclass[10pt]{article}
\usepackage[T1]{fontenc}
\usepackage{newtxtext,newtxmath}
\usepackage[margin=0.75in]{geometry}
\usepackage{amsmath}
\usepackage{booktabs,array}
\usepackage{xcolor}

\newcommand{\CPU}{\textsc{cpu}}
\newcommand{\GPU}{\textsc{gpu}}
\newcommand{\Hybrid}{\textsc{hybrid}}

\begin{document}

\section*{EdgeSched Phase Model}

EdgeSched profiles an executable \emph{phase}, not a model family.  A phase may
be an LLM decode, a TTS generator or codec, a diffusion job, or an auxiliary
stage.  Its static capability declaration is

\begin{equation}
\Phi_i = \left\langle
    d_i,\rho_i,g_i,\chi_i,
    \operatorname{role}_i,
    \operatorname{preempt}_i,
    \operatorname{queue}_i,
    \operatorname{deps}_i
\right\rangle .
\end{equation}

\begin{itemize}
    \item $d_i$: allowed devices, such as CPU-only, CPU-preferred,
          GPU-preferred, or GPU-required;
    \item $\rho_i$: fixed, whole-model migratable, or partitionable residency;
    \item $g_i$: validated split granularity (whole model, layer, component, or state);
    \item $\chi_i$: unmanaged GPU execution, cap-limited execution, or explicit
          SM partitioning;
    \item the remaining fields describe service role, preemption, queueing, and
          workflow dependencies.
\end{itemize}

For each supported configuration, offline profiling records a measured profile

\begin{equation}
\Psi_i(r,c,s,b\mid z)=\left\langle
    M_i^{G},M_i^{H},T_i^{\mathrm{first}},T_i^{\mathrm{steady}},
    Q_i,C_i^{\mathrm{transition}}
\right\rangle ,
\end{equation}

where $r$ is the residency configuration, $c$ the CPU-core budget, $s$ the
effective SM budget, $b$ the batch setting, and $z$ the input and runtime
context.  $M_i^G$ and $M_i^H$ are measured VRAM and host-memory demands; $Q_i$
contains workload-specific quality-of-service measurements such as TTFT, TTFA,
RTF, and slowdown.  Unsupported or unverified configurations are absent rather
than assigned optimistic estimates.

\section*{Current Workload Mapping}

\begin{center}
\begin{tabular}{>{\raggedright\arraybackslash}p{0.15\linewidth}
                >{\raggedright\arraybackslash}p{0.19\linewidth}
                >{\raggedright\arraybackslash}p{0.29\linewidth}
                >{\raggedright\arraybackslash}p{0.27\linewidth}}
\toprule
Phase & Service semantics & Allowed placement & Scheduling obligation \\
\midrule
LLM & resident, stateful, interactive & CPU, validated layer-hybrid, or full GPU
    & remain reachable; preserve KV/context across transitions \\
TTS & burst or session-warm, deadline-sensitive & CPU, validated component-hybrid,
    or full GPU & minimize first audio delay and meet subsequent playback deadlines \\
Diffusion & burst, queueable, GPU-required & GPU only
    & admit only when VRAM-safe; do not interrupt an admitted image job \\
Auxiliary stages & helper phases, normally bursty & measured CPU-preferred path;
    GPU only when profiling justifies it & avoid starving system, audio, and interactive threads \\
\bottomrule
\end{tabular}
\end{center}

These are workload instances of $\Phi_i$, not permanent classes tied to model
names.  In particular, the same TTS model may be a one-shot burst or remain warm
during an active voice session.

\section*{Runtime State and Decision Events}

At an event boundary, the scheduler observes

\begin{equation}
x_t=\left\langle
q_D,q_T,K,{\bigl\{r_i\bigr\}}_{i\in\mathcal A_t},M_G,S_G,U_C,v,h
\right\rangle ,
\end{equation}

including diffusion and TTS queues, LLM context/KV size, active residency
configurations, VRAM and SM state, CPU pressure, voice-session state, and recent
transition history.  Decisions occur on request arrival/completion, TTS queue
changes, voice-session changes, and backend failures; they are not recomputed
after every token or short utilization fluctuation.

The three principal operating states are

\begin{align*}
S_0 &: \text{no diffusion job: balance the resident LLM and optional TTS},\\
S_1 &: \text{diffusion running: protect an admitted GPU burst and interaction SLOs},\\
S_2 &: \text{diffusion queue nonempty: avoid restoring and immediately evicting again}.
\end{align*}

\section*{Hierarchical Feasibility: Capacity Before Compute Sharing}

The first decision selects one supported residency configuration $r_i$ for each
active phase.  A placement is capacity-feasible only if

\begin{align}
\sum_{i\in\mathcal A_t} M_i^G(r_i)+G_{\mathrm{safe}} &\le V_{\mathrm{GPU}},\\
\sum_{i\in\mathcal A_t} M_i^H(r_i)+R_H &\le V_{\mathrm{host}},\\
\sum_{i\in\mathcal A_t} c_i+R_C &\le C_{\mathrm{physical}}.
\end{align}

Only after these constraints pass does EdgeSched choose compute shares.  For
backends that provide verified disjoint partitions,

\begin{equation}
\sum_{i\in\mathcal G_t}s_i\le S_{\mathrm{GPU}}.
\end{equation}

For MPS active-thread caps, $s_i$ is recorded as a client limit rather than
claimed as an exclusive SM reservation.  Every shared configuration must also
satisfy measured interaction/audio constraints and a chosen diffusion slowdown
bound.  VRAM-compatible co-residency therefore does not imply execution-compatible
co-scheduling.

\section*{Joint Action and Cost}

The event-driven action is

\begin{equation}
a_t=\left({\bigl\{(r_i,c_i,s_i,b_i)\bigr\}}_{i\in\mathcal A_t},\mu\right),
\end{equation}

where $\mu$ is hold, urgent release, atomic restore, or reconfiguration.  For a
candidate action, the transition cost is measured as

\begin{equation}
C_{\mathrm{transition}}(a_t)=
C_{\mathrm{load}}+C_{\mathrm{state}}+C_{\mathrm{cutover}}+C_{\mathrm{recovery}}.
\end{equation}

The planner minimizes predicted end-to-end latency and interference rather than
an isolated task's execution time:

\begin{equation}
a_t^*=\arg\min_{a\in\mathcal F(x_t)}
\left[
T_{\mathrm{workflow}}(a\mid x_t)
+C_{\mathrm{transition}}(a)
+P_{\mathrm{SLO}}(a)
+P_{\mathrm{interference}}(a)
\right],
\end{equation}

where $\mathcal F(x_t)$ contains only profiled, capacity-safe configurations.
For a streaming LLM-to-TTS response, the first-audio objective includes

\begin{equation}
T_{\mathrm{first\ audio}}=
T_{\mathrm{transition}}
+T_{\mathrm{LLM}}^{\mathrm{first\ chunk}}
+T_{\mathrm{TTS}}^{\mathrm{first\ audio}},
\end{equation}

and every later audio chunk must be ready before the previous chunk finishes
playing.

\section*{Hysteresis and Safe Reconfiguration}

Let $a_{\mathrm{cur}}$ be the current action and $a'$ a feasible candidate.
EdgeSched changes configuration only when

\begin{equation}
J(a_{\mathrm{cur}}\mid x_t)-J(a'\mid x_t)
> C_{\mathrm{switch}}(a_{\mathrm{cur}},a')+\eta ,
\end{equation}

where $\eta$ is a hysteresis margin.  This prevents repeated model reloads or SM
reconfiguration in response to transient telemetry.  During a burst train,
EdgeSched may keep a degraded or hybrid residency configuration until the queue
drains instead of restoring after every job.

\end{document}
